\documentclass[10pt,twocolumn]{article}
\usepackage[a4paper,margin=18mm]{geometry}
\usepackage{kotex}
\usepackage{amsmath}
\usepackage{booktabs}
\usepackage{graphicx}
\usepackage{hyperref}

\title{한국어 도구 사용 에이전트의 채널별 압축 비용 전이 측정과 계약 보존 압축}
\author{Anonymous for review}
\date{}

\begin{document}
\maketitle

\begin{abstract}
본 연구는 다중 턴 도구 사용 에이전트에서 입력 압축의 명목상 토큰 절감이
추가 호출, 재시도, reasoning 및 출력 비용으로 얼마나 전이되는지를 측정한다.
입력을 지시·이력(I)과 도구 관측(T), 현재 생성을 출력(O)로 분리하고 한국어와
동일 의미 구조의 영어 태스크를 통제 비교한다.
총 880개 계획 run의 단계적 실험에서 균등 추출 압축은 표준화 비용을
줄였지만 성공률을 크게 낮췄다. 반면 의사결정 필드를 보존하는 탐색적
도구 계약 보존 압축(TCPC)은 80/80 성공을 유지하면서 비용을 12.61\% 절감했다.
\end{abstract}

\section{서론}
프롬프트 압축은 긴 입력의 추론 비용을 줄이지만\cite{jiang2023llmlingua,
jiang2024longllmlingua,pan2024llmlingua2}, 도구 사용 에이전트에서는 압축된
지시나 관측이 이후 행동 자체를 바꾼다. 입력 토큰만 비교하면 잘못된 도구 호출,
재시도, 조기 종료로 인한 비용과 유용성 변화를 놓칠 수 있다. 본 연구는
(1) I/T/O 채널별 비용 회계, (2) 원문을 복원한 반사실적 shadow-cost 기반 CDR,
(3) 한국어·영어 의미쌍 통제 실험, (4) 성공률을 먼저 통과시킨 뒤 비용을
최적화하는 배치 규칙을 제안한다. 추가로 구조적 도구 관측의 의사결정 필드를
보호하는 탐색적 Tool Contract-Preserving Compression(TCPC) 방법을 평가한다.

\section{관련 연구}
LLMLingua 계열은 중요 토큰 선택과 학습 기반 압축으로 일반 프롬프트의
효율과 정확도를 함께 평가한다\cite{jiang2023llmlingua,jiang2024longllmlingua,
pan2024llmlingua2}. Telegraph English는 고정 비율 삭제 대신 정보 밀도에
따른 구조적 재표현을 제안한다\cite{arbuzov2026telegraph}. 에이전트 문맥에서는
ACON이 관측과 이력을 함께 최적화하고\cite{acon2025}, OSWorld는 관측 및 이력
길이가 에이전트 성공에 영향을 줌을 보였다\cite{xie2024osworld}.
CAVEWOMAN은 입력·출력 압축이 비용과 정확도에 서로 다른 영향을 줄 수 있음을
보고한다\cite{cavewoman2026}. SkillZip은 실행 가능한 skill graph의 경계 계약을
보존해 재사용 문맥을 압축한다\cite{tan2026skillzip}. 본 연구는 이들과 달리
도구 관측 payload의 의사결정 필드를 보존하면서 다중 턴 trajectory의
비용을 I/T/O로 분리하고, 동일 행동 경로에 원문을 복원한 shadow-cost로 명목
절감과 행동 전이를 분해한다. 또한 한국어와 영어를 공통 의미 구조로 짝지어
비교하고, 토큰 절감만으로 선택한 조건과 성공률 우선 규칙의 결론이 어떻게
달라지는지 정량화한다.

\begin{table*}[t]
\centering
\caption{관련 연구와 평가 범위의 정성 비교. 보고 수치는 서로 다른 benchmark에서 측정되어 직접적인 우열 수치가 아니다.}
\resizebox{\textwidth}{!}{%
\begin{tabular}{lllll}
\toprule
연구 & 압축 대상 & 채널 구분 & 에이전트 실행 성공 & 행동 변화 비용 분해 \\
\midrule
LLMLingua 계열 & 일반 프롬프트 & 입력 중심 & 아니오 & 아니오 \\
CAVEWOMAN & 언어 입력·출력 & I/O & 아니오 & 응답 비용 비교 \\
ACON & 장기 에이전트 관측·이력 & 관측/이력 & 예 & peak token 중심 \\
SkillZip & 에이전트 skill graph & 실행 계약 & 예 & token/검색 성능 \\
본 연구 & 다중 턴 도구 trajectory & I/T/O & 예 & shadow-cost와 CDR \\
\bottomrule
\end{tabular}}
\end{table*}

\section{방법}
각 요청의 콘텐츠 입력을 I와 T로 분류하고 현재 모델 생성을 O로 분류한다.
로컬 채널은 의미 필드의 이름과 값만 세며 JSON 구두점과 전송 전용
\texttt{type/role/id/call\_id/status/encrypted\_content} 래퍼는 제외한다.
provider 직렬화 차이는
어느 채널에도 임의 배분하지 않고 residual로 둔다.
단, 도구 결과는 adapter가 canonical JSON 문자열로 보내므로 T는 그 실제
직렬화 문자열(키와 구두점 포함)을 토큰화한다.
압축 조건은 B0, I-H, T-H, O이며 목표 유지율은 0.4이다.
GPT-5.6 Terra는 temperature 파라미터를 지원하지 않아 해당 필드를 생략하고
모델 기본값을 사용한다. 이는 성공한 스모크 응답을 관찰하기 전에 고정했다.
실제 과금 비용은 입력을 비캐시, 캐시 읽기, 캐시 쓰기로 분리해 계산하고,
list-price 비용은 모든 입력을 비캐시 정가로 표준화한다. 272K 입력 초과 요청은
요청 전체에 장문 요율을 적용하되 주 분석 요청은 240K 미만으로 제한한다.

\begin{align}
G_i &= C_{\mathrm{shadow},i} - C_{\mathrm{compressed},i},\\
B_i &= C_{\mathrm{shadow},i} - C_{\mathrm{baseline},i} + K_i,\\
\widehat{\mathrm{CDR}} &= \frac{\sum_i B_i}{\sum_i G_i}.
\end{align}

한국어·영어 surface는 40개 공통 의미 구조에서 생성하며 각 의미쌍을
bootstrap 재표집 단위로 사용한다. B0는 인스턴스별 3회 반복한다.
최초 실험은 B0 240회와 I-H/T-H/O 각 80회로 구성한다. 최초 결과를 바탕으로,
후속 실행 전에 I-H와 T-H의 유지율 0.6·0.8 각 80회를 동결해 320회 추적
실험을 수행했다.
성공률 하락의 단측 95\% 상한이 5\%p 이하이고 비용 절감의 pair-bootstrap
95\% 신뢰구간 하한이 0보다 큰 경우만 안전한 절감으로 판정한다. 네 셀의
성공률 차이 검정은 Holm 보정을 적용한다.

\subsection{탐색적 TCPC}
비율 추적 결과 이후 별도의 탐색 단계에서 TCPC를 설계했다.
도구 adapter가 의사결정에 쓰이지 않는 필드를 선언하면 해당 필드만 제거하고,
ID, 상태, 금액, 날짜, evidence code 등 나머지 계약 필드는 원문 그대로 둔다.
보호 필드를 유지한 표현이 목표 유지율보다 길어도 추가 삭제하지 않는
안전 바닥(safety floor)을 사용한다. 본 태스크에서는 \texttt{notes}만
비계약 필드로 선언했으며, T-H 유지율 0.8에서 80회를 실행했다. 이 단계는 앞선
결과로부터 도출되었으므로 확증 결과가 아니라 탐색 결과로 구분한다.

\section{결과}
480개 계획 ordinal을 모두 실행했고, 회계 불변식을 벗어난 12개 I-H run을
사전 규칙에 따라 제외하여 468개를 주 분석에 사용했다. B0와 O는 모두
성공했지만 T-H는 62.5\%, I-H는 8.8\%만 성공했다. 제외된 I-H를 실패로 보는
민감도 집계에서는 I-H 성공률이 7.5\%(6/80)다.

평균 list-price 비용은 B0 대비 I-H에서 48.7\%, T-H에서 37.1\%, O에서
1.4\% 감소했다. 그러나 비용 감소는 품질 보존을 뜻하지 않는다. I-H의
낮은 비용은 주로 조기 종료와 적은 도구 호출을 동반했다.

\begin{table*}[t]
\centering
\caption{조건별 유효 run 요약. 비용은 trajectory당 USD, reasoning은 provider token 평균이다.}
\resizebox{\textwidth}{!}{%
\begin{tabular}{lrrrrrrrrr}
\toprule
조건 & $n$ & 성공률 & List 비용 & Billed 비용 & Reasoning & 턴 & Tool call & 실현률 & CDR \\
\midrule
B0  & 240 & 1.000 & 0.007308 & 0.005387 & 11.2  & 2.50 & 1.50 & 1.000 & -- \\
I-H &  68 & 0.088 & 0.003751 & 0.003481 & 129.2 & 1.38 & 0.38 & 0.339 & -19.563 \\
T-H &  80 & 0.625 & 0.004595 & 0.004609 & 90.9  & 2.14 & 1.14 & 0.388 & -0.629 \\
O   &  80 & 1.000 & 0.007203 & 0.005621 & 12.1  & 2.50 & 1.50 & 1.000 & -- \\
\bottomrule
\end{tabular}
}
\end{table*}

Aggregate CDR은 I-H에서 $-19.56$ (pair bootstrap 95\% CI
$[-29.90,-11.91]$), T-H에서 $-0.629$ (95\% CI $[-0.983,-0.368]$)였다.
두 값이 음수라는 것은 추가 행동 비용이 절감을 상쇄하지 않았다는 뜻이지만,
본 실험에서는 압축이 오히려 행동을 중단시키거나 잘못된 인자를 생성해 비용을
더 낮춘 경우가 많았다. 성공률 저하와 함께 해석하지 않으면 잘못된 최적화
결론에 이를 수 있다. McNemar exact test도 B0 대비 I-H와 T-H의 성공 저하를
각각 $p<10^{-18}$, $p<10^{-8}$로 나타냈다.

비용 전이 분해에서 reasoning 증가는 I-H와 T-H에 각각 \$0.0961,
\$0.0765를 더했지만, 도구 관측 입력 변화가 각각 -\$0.1858,
-\$0.1172였고 다른 출력·입력 감소까지 합쳐 순 행동 비용은 음수가 됐다.
실패 유형은 I-H의 조기 종료 52건과 인자 오류 10건, T-H의 인자 오류
30건이었다. 탐색적 언어별 성공률은 I-H에서 한국어 0\%(0/34), 영어
17.6\%(6/34)였고 T-H는 두 언어 모두 62.5\%였다.

\begin{table}[t]
\centering
\caption{유지율 추적 및 구조 보존 압축 결과. 절감률은 B0 대비 list-price 비용이다.}
\resizebox{\columnwidth}{!}{%
\begin{tabular}{lrrrr}
\toprule
방법 & 유지율 & 계획 $n$ & 성공률 & 비용 절감 \\
\midrule
I-H 균등 추출 & 0.6 & 80 & 12.5\% & 40.26\% \\
I-H 균등 추출 & 0.8 & 80 & 46.3\% & 2.21\% \\
T-H 균등 추출 & 0.6 & 80 & 78.8\% & 26.43\% \\
T-H 균등 추출 & 0.8 & 80 & 82.5\% & 14.71\% \\
T-H TCPC & 0.8 & 80 & 100.0\% & 12.61\% \\
\bottomrule
\end{tabular}}
\end{table}

어느 균등 추출 셀도 5\%p 성공률 비열화 기준을 통과하지 못했다. 반면
TCPC는 80/80 성공했고 성공 손실의 단측 95\% 상한은 3.68\%였다.
비용 절감의 95\% CI는 [10.67\%, 13.94\%]였으며 안전성과 양의 절감을 모두
통과했다. 동일 태스크의 T-H 균등 추출 0.8과 직접 비교하면 TCPC만
성공한 경우가 14개, 반대가 0개였다(McNemar exact $p=0.000122$).

% 생성 SVG: results/figures/figure1_channel_saving.svg,
% results/figures/figure2_cdr_components.svg. 제출본에서는 PDF로 변환한다.

\section{논의 및 한계}
결과는 채널별 압축의 비용 효과와 유용성 효과가 분리되어야 함을 보여준다.
특히 I-H 압축은 reasoning token 평균을 높이면서도 조기 종료가 많았고,
T-H 압축은 비용을 줄였지만 인자 오류와 37.5\%p 성공률 하락을 만들었다.
반면 O 조건은 B0와 같은 성공률을 유지했으나 list-price 감소는 1.4\%로 작았다.
따라서 제안 배치 규칙은 토큰 또는 비용 절감률이 가장 큰 I-H를 선택하지 않고,
성공률 문턱을 통과한 O 또는 TCPC만 후보로 남긴다. 이는 동일 실행
자료에서 토큰 중심 평가와 성공률 우선 평가가 서로 다른 결론을 낸다는
정량적 비교다.

합성 태스크, 단일 모델, 비공개 tokenizer/직렬화 residual, 자동 caching,
언어별 surface 자연스러움과 모델 기본 sampling의 실행 비결정성이 한계다.
또한 회계-invalid 12건이 모두 I-H에 집중되어 결측이 무작위가 아니므로,
I-H의 유효-run 추정치는 보수적으로 해석하고 계획 run 기준 민감도 결과를
함께 보고한다. 한·영 차이와 중립/민감 하위집단 결과는 다중비교 교정 없는
기술통계이므로 확증적 언어 효과로 해석하지 않는다.
TCPC 역시 합성된 구조적 관측과 사전에 선언한 단일 비계약 필드에서만
평가했으며, 비정형 관측 또는 다른 도구 schema로 일반화되었다고 볼 수 없다.
후속 연구에서는 여러 모델·실환경 benchmark·압축기를 사용하고, 계약 필드
선언 오류와 compressor 자체 비용까지 포함해 검증해야 한다.

\section{결론}
도구 사용 에이전트의 압축은 입력 토큰 감소만으로 평가할 수 없다. 채널별
shadow-cost와 성공률을 함께 사용하면 실패로 싸진 실행을 효율 향상으로
오인하지 않는다. 균등 추출 압축은 유지율 0.8에서도 안전하지 않았지만,
탐색적 TCPC는 의사결정 계약을 보존하여 성공률을 유지하면서 유의한
비용 절감을 달성했다.

\bibliographystyle{plain}
\bibliography{references}
\end{document}
